% Pursuing Stacks (cotes 134-2 to 134-8): archive cote 134-2, site page 58 (PDF page 59).
% appendix (letters to Breen). Draft compared: ps1.tex:1404-1424.
% AI-transcribed by Claude Opus 5.5 (claude-opus-5-5, adaptive thinking, effort high), Message Batches API, batch ps1 (msgbatch_01Rv5hCaNTtUt4JHx1yqoHPL), 2026-11-12; not yet checked by eye.
\scanpage{134-2/58}
\mspage{-13 -}
operates on $L_1$ leaving $\pi_1$ invariant, and it operates \emph{via} $\pi_0$ . The
principal cohomological invariant of this situation is evidently the Postnikoff-%
Sinh invariant
\[ \alpha \in H^3(\pi_0, \pi_1) . \]

I have met these animals - without even looking for them - in many
situations, which I will not list now (I came across them recently \emph{a propos}
the classification of ``ordinary'' formal groups over a perfect field, in terms
of \emph{affine} algebraic groups, and \emph{commutative} formal groups, related by the
strict Gr-structures of this type (except that one has to use this formalism
in an arbitrary topos (not merely in \deleted{(\unsure{Ens})}\inserted{(Sets)}) - to make explicit the yoga that
``the transcendent character of a formal group is concentrated essentially in
the commutative formal group\hand{s}'', discovered it seems by Dieudonné ...). The
question which I wish to raise is the generalisation to groups in (n-cat),
where I expect to find a non-commutative chain complex
\[ L_n \to L_{n-1} \to \dots \to L_1 \to L_0 \to 1 \]
with supplementary structures doubtless of the type of $\theta$ , but what are they?
It is understood that the topological significance of \deleted{the}\inserted{such} structures is that
they express exactly the ``truncated homotopy type in dim $\leq$ n'' of topological
groups, \deleted{and even}\inserted{or equivalently} the homotopy type in dim $\leq n + 1$ of \deleted{connected}\inserted{pointed connected} topological
spaces ...). Have you candidates to propose?
\margin{espacer}
* * *

\inserted{(App.10)} Your reflections on biduality and homology, \deleted{are attached to}\inserted{however formal, tie in with} a crowd of
developments, of which \inserted{only} some \deleted{only} exist at present, and others would demand
\deleted{some} considerable \deleted{development}\inserted{work still}. Here are the reminiscences which your naive
questions bring to mind:
(A) The formalism of the $Rf_!$ , $Rf^!$ (combined with $Rf_*$ , $Lf^*$ , $\overset{\hand{$\mathbb{L}$}}{\otimes}$ and \hand{$\mathbb{R}$}Hom,
``the six operations'') carries implicitly in itself the definition of homology
and the essential identit\over{i}{y}\deleted{es that underly a}\inserted{between} homology and cohomology. One now
has this formalism for quasi-coherent sheaves on schemes - seminar Hartshorne
(Springer L.N. 20) - for the topological spaces and arbitrary sheaves of
\mspage{\hand{58}}
